\subsection{\texorpdfstring{Bounded measures and linear forms on \(\mathcal{C}^{b}(T)\)}{Bounded measures and linear forms on \textbackslash mathcal\{C\}\^{}\{b\}(T)}}

PROPOSITION 5. --- Let T be a completely regular space, and I a continuous complex linear form on the normed space \(\mathcal{C}^{b}(T;C)\) . In order that there exist a bounded complex measure \(\theta\) on T such that \(\theta(f)=\mathrm{I}(f)\) for all \(f\in\mathcal{C}^{b}(T;C)\) , it is necessary and sufficient that the following condition be satisfied:

\begin{enumerate}
\def\labelenumi{(\Alph{enumi})}
\setcounter{enumi}{12}
\tightlist
\item
  For every number \(\varepsilon > 0\) , there exists a compact subset \(K\) of \(T\) such that the relations \(g \in \mathcal{C}^b(T; \mathbf{C})\) , \(|g| \leqslant 1\) , \(g_K = 0\) imply \(|\mathrm{I}(g)| \leqslant \varepsilon\) .
\end{enumerate}

The measure \(\theta\) is then unique.

Uniqueness follows from Prop. 2 of No.~1. Let us show that the condition (M) is necessary. Let \(\theta\) be a bounded complex measure; let K be a compact set such that \(|\theta|^{\bullet}(\mathrm{T}-\mathrm{K}) \leqslant \varepsilon\) (§1, No.~2, Remark 3). The hypotheses \(|g| \leqslant 1\) , \(g_{\mathbf{K}} = 0\) imply \(|g| \leqslant \varphi_{\mathbf{C}_{\mathbf{K}}}\) , therefore \(|\theta(g)| \leqslant |\theta|^{\bullet}(\varphi_{\mathbf{C}_{\mathbf{K}}}) \leqslant \varepsilon\) .

Let us pass to the proof of sufficiency. Let X be the Stone--Čech compactification of T (GT, IX, §1, Exer. 7; or TG, IX, §1, No.~6). For every function \(f \in \mathcal{C}(X; \mathbf{C})\) , set \(\nu(f) = I(f_{\mathrm{T}})\) ; we define in this way a continuous linear form on \(\mathcal{C}(X; \mathbf{C})\) , that is, a complex measure on the compact space X. Let \(\varepsilon\) be a number \(>0\) , K a compact set satisfying (M); the function \(\varphi_{\mathbf{C}_{\mathrm{K}}}\) being lower semi-continuous and positive on X, the formula (2) gives us the following relations, where \(\mathcal{G}\) denotes the set of functions \(g \in \mathcal{C}(X; \mathbf{C})\) such that \(|g| \leqslant \varphi_{\mathbf{C}_{\mathrm{K}}}\) :

\[
| \nu | ^ {\bullet} (\mathrm{X} - \mathrm{K}) = \sup _ {g \in \mathcal {G}} | \nu (g) | = \sup _ {g \in \mathcal {G}} | \mathrm{I} (g _ {\mathrm{T}}) | \leqslant \varepsilon .
\]

Let \((\mathrm{K}_n)_{n\geqslant 1}\) be a sequence of compact subsets of T, such that each \(\mathbf{K}_n\) satisfies (M) for \(\varepsilon = 1 / n\) , and let \(\mathrm{S} = \bigcup_{n}\mathrm{K}_{n}\) ; S is a Borel set in X, contained in T, and \(|\nu|^{\bullet}(\mathrm{X} - \mathrm{T})\leqslant |\nu |^{\bullet}(\mathrm{X} - \mathrm{S})\leqslant |\nu |^{\bullet}(\mathrm{X} - \mathrm{K}_n)\leqslant 1 / n\) for all \(n\) , so that T is \(\nu\) -measurable and \(\nu\) is concentrated on T. Let \(f\) be a bounded continuous function on T; since X is the Stone-Čech compactification of T, \(f\) may be extended by continuity to a function \(g\in \mathcal{C}(\mathrm{X};\mathbf{C})\) . Now let \(\mu\) be the measure induced by \(\nu\) on T; one has \(\mu(f) = \nu(f^0)\) . Since \(\nu\) is concentrated on T, the functions \(f^0\) and \(g\) are equal \(\nu\) -almost everywhere, therefore \(\mu(f) = \nu(g) = I(g_T) = I(f)\) , which completes the proof.

COROLLARY. --- With notations as in Prop. 5, suppose that there exists a bounded positive measure \(\mu\) on T such that \(|\mathrm{I}(f)| \leqslant \mu(|f|)\) for all \(f \in \mathcal{C}^b(\mathrm{T};\mathbf{C})\) ; then there exists a complex measure \(\theta\) on T such that \(\theta(f) = \mathrm{I}(f)\) for all \(f \in \mathcal{C}^b(\mathrm{T};\mathbf{C})\) .

